\subsection{Compact linear mappings}

DEFINITION 1. --- Let E be a topological vector space and F a separated topological vector space. A linear mapping u from E into F is said to be compact if there exists a neighbourhood V of 0 in E such that \(u(V)\) is a relatively compact subset of F.

We denote by \(\mathcal{L}^{\mathrm{c}}(\mathrm{E};\mathrm{F})\) the set of compact linear mappings from E into F; we also denote by \(\mathcal{L}^{\mathrm{c}}(\mathrm{F})\) the set \(\mathcal{L}^{\mathrm{c}}(\mathrm{F};\mathrm{F})\) .

Remarks. --- 1) A neighbourhood of 0 in E absorbs every bounded subset of E; the image of a bounded subset of E under a compact linear mapping from E into F is therefore a relatively compact subset of F.

\begin{enumerate}
\def\labelenumi{\arabic{enumi})}
\setcounter{enumi}{1}
\item
  Let E be a semi-normed space and let B be the unit ball of E. For a linear mapping u from E into a separated topological vector space F to be compact, it is necessary and sufficient that \(u(\mathrm{B})\) is a relatively compact subset of F.
\item
  Let F be a separated topological vector space. For a vector subspace E of F to be finite-dimensional, it is necessary and sufficient that the canonical injection of E into F be compact (EVT, I, p. 15, th. 3).
\item
  For every topological vector space E over the field R, we denote by \(\mathrm{E}_{(\mathbf{C})}\) the complexified topological vector space of E (EVT, II, p. 65). We identify E with a real topological vector subspace of \(\mathrm{E}_{(\mathbf{C})}\) by the mapping \(x \mapsto 1 \otimes x\) ; the underlying real topological vector space of \(\mathrm{E}_{(\mathbf{C})}\) is then the topological direct sum of E and iE.
\end{enumerate}

Let \(u\) be a linear mapping from E into a separated topological vector space F over \(\mathbf{R}\) . Denote by \(u_{(\mathbf{C})}\) the \(\mathbf{C}\)-linear mapping from \(\mathrm{E}_{(\mathbf{C})}\) into \(\mathrm{F}_{(\mathbf{C})}\) which extends \(u\) . The sets of the form \(V + iV\) , where \(V\) is a neighbourhood of 0 in E, form a fundamental system of neighbourhoods of 0 in \(\mathrm{E}_{(\mathbf{C})}\) . For the set \(u_{(\mathbf{C})}(V + iV) = u(V) + iu(V)\) to be relatively compact in \(\mathrm{F}_{(\mathbf{C})}\) , it is necessary and sufficient that \(u(V)\) be relatively compact in F. Consequently, for the \(\mathbf{C}\)-linear mapping \(u_{(\mathbf{C})}\) to be compact, it is necessary and sufficient that the \(\mathbf{R}\)-linear mapping \(u\) be so.

In the same vein, when \(v\) is a continuous C-linear mapping from \(\mathrm{E}_{(\mathbf{C})}\) into a C-topological vector space G, for \(v\) to be compact it is necessary and sufficient that its restriction to E be compact.

\begin{enumerate}
\def\labelenumi{\arabic{enumi})}
\setcounter{enumi}{4}
\item
  Let E be a topological vector space, and let F be a separated topological vector space. Let \(F_{1}\) be a closed vector subspace of F and let j be the canonical injection of \(F_{1}\) into F. For a linear map u from E into \(F_{1}\) to be compact, it is necessary and sufficient that the linear map \(j \circ u\) from E to F be so.
\item
  Let \((\mathrm{G}_{i})_{i\in\mathrm{I}}\) be a family of separated topological vector spaces, and for every \(i\in I\) , let \(u_{i}\) be a compact linear map from a topological vector space E into \(G_{i}\) . If the set I is finite, or if E is a semi-normed space, the linear map \(x\mapsto(u_{i}(x))_{i\in\mathrm{I}}\) from E into \(\prod_{i\in I}G_{i}\) is compact (Remark 2 and TG, I, p. 63, Th. 3).
\end{enumerate}

PROPOSITION 1. --- Let E be a topological vector space and F a separated topological vector space.

\begin{enumerate}
\def\labelenumi{\alph{enumi})}
\item
  Every compact linear map from E into F is continuous.
\item
  Every continuous linear map of finite rank from E into F is compact;
\item
  The set \(\mathcal{L}^{\mathrm{c}}(\mathrm{E};\mathrm{F})\) of compact linear mappings from E into F is a vector subspace of \(\mathcal{L}(\mathrm{E};\mathrm{F})\) .
\end{enumerate}

Let us prove a). Let \(v: E \rightarrow F\) be a compact linear mapping and U a neighborhood of 0 in F. Choose a neighborhood V of 0 in E such that \(v(V)\) is relatively compact in F. Then \(v(V)\) is bounded, and therefore there exists a real number \(\lambda > 0\) such that the set \(v(\lambda V)\) is contained in U; consequently, v is continuous.

Let \(u \colon \mathrm{E} \to \mathrm{F}\) be a continuous linear map of finite rank. Since the image \(u(\mathrm{E})\) of \(u\) is separated and finite-dimensional, there exists a compact neighborhood A of 0 in \(u(\mathrm{E})\) . The set \(\mathrm{V} = u^{-1}(\mathrm{A})\) is a neighborhood of 0 in E, and we have \(u(\mathrm{V}) \subset \mathrm{A}\) , so the linear map \(u\) is compact. This proves the assertion \(b\) ).

Let \(u_{1}\) and \(u_{2}\) be compact linear mappings from E into F. Let \(V_{1}\) and \(V_{2}\) be neighborhoods of 0 in E such that the sets \(u_{1}(V_{1})\) and \(u_{2}(V_{2})\) are relatively compact in F. Set \(V = V_{1} \cap V_{2}\) and \(u = u_{1} + u_{2}\) . Then \(u(V)\) is contained in \(u_{1}(V_{1}) + u_{2}(V_{2})\) and therefore is relatively compact in F, and the linear map u is compact. It is immediate that \(\mathcal{L}^{c}(E;F)\) is stable under homothety. This proves that \(\mathcal{L}^{c}(E;F)\) is a vector subspace of \(\mathcal{L}(E;F)\) , whence assertion c).

PROPOSITION 2. --- Let E be a semi-normed space and F a separated topological vector space. Suppose that every closed bounded subset of F is complete. Then the vector subspace \(\mathcal{L}^{\mathrm{c}}(\mathrm{E};\mathrm{F})\) of \(\mathcal{L}(\mathrm{E};\mathrm{F})\) is closed for the topology of bounded convergence.

Let \(u\) be an element of \(\mathcal{L}(\mathrm{E};\mathrm{F})\) adherent to \(\mathcal{L}^{\mathrm{c}}(\mathrm{E};\mathrm{F})\) for the topology of bounded convergence. Let B denote the unit ball of E and V a neighborhood of 0 in F. Choose a neighborhood W of 0 in F such that \(\mathrm{W} + \mathrm{W} \subset \mathrm{V}\) . By hypothesis there exists an element \(v\) of \(\mathcal{L}^{\mathrm{c}}(\mathrm{E};\mathrm{F})\) such that \((u - v)(\mathrm{B}) \subset \mathrm{W}\) . The set \(v(\mathrm{B})\) being relatively compact in F, there exists a finite subset M of F such that \(v(\mathrm{B}) \subset \mathrm{M} + \mathrm{W}\) , whence \(u(\mathrm{B}) \subset \mathrm{M} + \mathrm{V}\) . Since this holds for every V, the set \(u(\mathrm{B})\) is precompact in F, hence so is its closure (TG, II, prop. 1, p. 30). Since the latter is closed and bounded, it is complete by the hypothesis made on F, and consequently is compact (TG, II, cor., p. 30). It follows that \(u(\mathrm{B})\) is relatively compact. This proves that \(u\) belongs to \(\mathcal{L}^{\mathrm{c}}(\mathrm{E};\mathrm{F})\) .

If F is locally convex, the hypothesis of the proposition means that F is quasi-complete (EVT, III, p. 8, def. 6).

COROLLARY. --- Let E and F be Banach spaces. The set \(\mathcal{L}^{\mathrm{c}}(\mathrm{E};\mathrm{F})\) is a closed vector subspace of the Banach space \(\mathcal{L}(\mathrm{E};\mathrm{F})\) . It contains the closure of the space \(\mathcal{L}^{\mathrm{f}}(\mathrm{E};\mathrm{F})\) of continuous linear mappings of finite rank from E into F.

This follows from Propositions 1 and 2.

\pandocbounded{\includegraphics[keepaspectratio]{images/part001_746f5095.jpg}}

There exist Banach spaces E and F such that \(\mathcal{L}^{\mathrm{f}}(\mathrm{E};\mathrm{F})\) is not dense in \(\mathcal{L}^{\mathrm{c}}(\mathrm{E};\mathrm{F})\) (cf. Remark 6 of III, p. 16 and Theorem 4 of III, p. 19, b)).

PROPOSITION 3. --- Let \(\mathrm{E}_{1}\) , E, F, \(\mathrm{F}_{1}\) be topological vector spaces, with F and \(\mathrm{F}_{1}\) assumed separated. Let v: \(\mathrm{E}_{1}\) → E, u: E → F, and w: F → \(\mathrm{F}_{1}\) be linear maps. If v and w are continuous and u is compact, then w ∘ u ∘ v is compact.

By hypothesis, there exists a neighborhood V of 0 in E such that \(u(\mathrm{V})\) is relatively compact in F. Set \(\mathrm{U} = v^{-1}(\mathrm{V})\) . Then U is a neighborhood of 0 in \(\mathrm{E}_{1}\) and its image under \(w \circ u \circ v\) is contained in \(w(u(\mathrm{V}))\) , hence is relatively compact in \(\mathrm{F}_{1}\) . Consequently, the linear map \(w \circ u \circ v\) is compact.

Let E be a separated topological vector space; by props. 1 and 3, \(\mathcal{L}^{\mathrm{c}}(\mathrm{E})\) is a two-sided ideal of the algebra \(\mathcal{L}(\mathrm{E})\) . When E is a Banach space, \(\mathcal{L}(\mathrm{E})\) is a Banach algebra and \(\mathcal{L}^{\mathrm{c}}(\mathrm{E})\) is a closed two-sided ideal of \(\mathcal{L}(\mathrm{E})\) (cor. of prop. 2). It is a proper ideal if E is of infinite dimension (cf. remark 3, p. 2).

COROLLARY. — Let E be a Hilbert space. The space \(\mathcal{L}^{\mathrm{c}}(\mathrm{E})\) is a closed self-adjoint two-sided ideal of \(\mathcal{L}(\mathrm{E})\) . In particular, it is a star algebra.

Indeed, the space \(\mathcal{L}^{c}(E)\) is a closed two-sided ideal, hence self-adjoint, of \(\mathcal{L}(E)\) (lemma 15 of I, p. 122).

PROPOSITION 4. — Let E be a topological vector space, \(\widehat{\mathrm{E}}\) its separated completion and \(j\) the canonical map from E into \(\widehat{\mathrm{E}}\) . Let \(u\) be a compact linear map from E into a separated topological vector space F. There exists a unique compact linear map \(v\) from \(\widehat{\mathrm{E}}\) into F such that \(u = v \circ j\) .

Identify F with a topological vector subspace of \(\widehat{\mathrm{F}}\) and let us denote \(\widehat{u}\colon \widehat{\mathrm{E}}\to \widehat{\mathrm{F}}\) the unique continuous linear map such that \(\widehat{u}\circ j\) coincides with \(u\) in E. Since \(u\) is compact, there exists a neighborhood V of 0 in E and a compact subset A of F such that \(u(\mathrm{V})\subset \mathrm{A}\) . We have \(\widehat{u}(j(\mathrm{V}))\subset \mathrm{A}\) , whence \(\widehat{u} (\overline{j(\mathrm{V})})\subset \mathrm{A}\) . Now \(\overline{j(\mathrm{V})}\) is a neighborhood of 0 in \(\widehat{\mathrm{E}}\) (TG, III, p. 24, prop. 7). Consequently, the image of \(\widehat{u}\) is contained in F and the continuous linear map \(v\colon \widehat{\mathrm{E}}\rightarrow \mathrm{F}\) derived from \(\widehat{u}\) by passing to the subspace is compact. Since \(j(\mathrm{E})\) is dense in \(\widehat{\mathrm{E}}\) , the map \(v\) is the unique continuous linear map from \(\widehat{\mathrm{E}}\) into F such that \(u = v\circ j\) .

PROPOSITION 5. --- Let E be a locally convex space, F a separated topological vector space, and u a compact linear map from E into F. There exist a Banach space G, a continuous linear map v from E into G, and a compact linear map w from G into F such that \(u = w \circ v\) .

Let V be a neighborhood of 0 in E such that \(u(\mathrm{V})\) is relatively compact in F. Let \(p\) be a continuous semi-norm on E such that V contains the unit ball of \(p\) (EVT, II, p. 26, cor.). Let us denote by \(\mathrm{E}_p\) the semi-normed space obtained by endowing E with the semi-norm \(p\) . The map \(u\) is a compact linear map from \(\mathrm{E}_p\) into F. The separated completion G of \(\mathrm{E}_p\) is a Banach space. Denote by \(v\) the canonical linear map from \(\mathrm{E}_p\) into G. By prop. 4, there exists a compact linear map \(w: \mathrm{G} \to \mathrm{F}\) such that \(u = w \circ v\) . On the other hand, v is a continuous map from E into G because the topology of E is finer than that of \(\mathrm{E}_p\) .

